\chapter{Stopień, indeks punktu stałego i twierdzenie Lefschetza}
\label{ch:lefschetz-v10}

\section{Dlaczego ślad działania na homologii wykrywa punkty stałe?}
Niech $X$ będzie skończonym kompleksem CW, a $f:X\to X$
odwzorowaniem ciągłym. Zwykły stopień mierzy działanie
$f$ na najwyższej homologii sfery. Dla przestrzeni, które
mają niezerową homologię w wielu stopniach, właściwym
uogólnieniem jest naprzemienna suma śladów działań
$f_*$ na wszystkich grupach homologii. Wynik może wymusić
istnienie punktu stałego, chociaż nie daje jego współrzędnych.
Rozdział oddziela twierdzenie dla dowolnych skończonych
kompleksów CW od dokładniejszego wzoru indeksowego na
gładkich rozmaitościach.

\section{Liczba Lefschetza i formuła śladu na łańcuchach}
\begin{definicja}[Liczba Lefschetza]
Niech $X$ będzie skończonym kompleksem CW, a $f:X\to X$
odwzorowaniem ciągłym. Ponieważ $H_k(X;\Q)$ są
skończeniewymiarowe i znikają dla $k>\dim X$, definiujemy
\[
 L(f)=\sum_{k\ge0}(-1)^k\operatorname{tr}
 \bigl(f_*:H_k(X;\Q)\to H_k(X;\Q)\bigr)\in\Z.
\]
Całkowitość wynika z tego, że $f_*$ zachowuje obrazy
wolnych części grup $H_k(X;\Z)$; macierze całkowite
mają całkowite ślady. Liczba $L(f)$ nie zależy od
wyboru bazy homologii.
\end{definicja}

\begin{lemat}[Lemat śladu Hopfa]
\label{lem:hopf-trace-v10}
Niech $(C_*,\partial)$ będzie ograniczonym kompleksem
skończeniewymiarowych przestrzeni wektorowych nad ciałem
$K$, a $F_*:C_*\to C_*$ odwzorowaniem łańcuchowym.
Wówczas
\[
 \sum_k(-1)^k\operatorname{tr}(F_k\mid C_k)
 =\sum_k(-1)^k\operatorname{tr}(H_k(F)\mid H_k(C)).
\]
\end{lemat}
\begin{proof}
Niech $Z_k=\ker\partial_k$ oraz $B_k=\operatorname{im}\partial_{k+1}$.
Wszystkie podprzestrzenie są niezmiennicze względem $F$.
Dokładne ciągi
$0\to Z_k\to C_k\to B_{k-1}\to0$ oraz
$0\to B_k\to Z_k\to H_k(C)\to0$
rozszczepiają się jako ciągi przestrzeni wektorowych.
Dla endomorfizmu zachowującego podprzestrzeń ślad jest
sumą śladu ograniczenia i śladu odwzorowania indukowanego
na ilorazie. Otrzymujemy
$\operatorname{tr}(F_k\mid C_k)=
 \operatorname{tr}(H_k(F))+
 \operatorname{tr}(F_k\mid B_k)+
 \operatorname{tr}(F_{k-1}\mid B_{k-1})$.
W naprzemiennej sumie dwa ostatnie wyrazy znoszą się
po przesunięciu indeksu. To daje tezę.
\end{proof}

\begin{wniosek}[Homotopie, deformacje i charakterystyka Eulera]
Jeśli $f,g:X\to X$ są homotopijne, to $L(f)=L(g)$.
Ponadto $L(\id_X)=\chi(X)$, gdzie
$\chi(X)=\sum_k(-1)^k\dim_\Q H_k(X;\Q)$.
Jeżeli $X$ jest niepusty i kontraktylny, to $L(f)=1$
dla każdego ciągłego $f:X\to X$.
\end{wniosek}
\begin{proof}
Homotopijne odwzorowania indukują te same homomorfizmy
w homologii. Dla $\id_X$ wszystkie ślady są wymiarami.
Dla przestrzeni kontraktylnej jedyną niezerową grupą
homologii racjonalnej jest $H_0(X;\Q)\cong\Q$,
a ciągłe odwzorowanie przestrzeni spójnej działa na niej identycznie.
\end{proof}

\section{Twierdzenie Lefschetza dla skończonych kompleksów}
\begin{twierdzenie}[Lefschetz — istnienie punktu stałego]
\label{thm:lefschetz-v10}
Jeśli $X$ jest skończonym kompleksem CW, $f:X\to X$
ciągłym odwzorowaniem i $L(f)\ne0$, to istnieje
$x\in X$ takie, że $f(x)=x$.
\end{twierdzenie}
\begin{proof}
Najpierw rozważmy skończony wielokompleks $X=|K|$.
Zapiszemy dowód kontrapozycji, a następnie wyprowadzimy
wersję dla dowolnego skończonego CW-kompleksu przez
osadzenie jako retraktu otoczenia w przestrzeni euklidesowej.

Załóżmy zatem dla sprzeczności, że $X=|K|$ i $f$
nie ma punktów stałych. Umieśćmy $|K|$ w przestrzeni
euklidesowej. Dzięki zwartości
$\delta=\min_{x\in X}\|f(x)-x\|>0$.
Przez dostatecznie drobny podział barycentryczny
kompleksu docelowego i wyjściowego możemy założyć,
że każdy sympleks $\sigma$ tego samego drobnego podziału
ma średnicę mniejszą niż $\delta/4$ (jednostajna
ciągłość $f$ nie jest tu potrzebna do oszacowania odległości
$f(\sigma)$ od $\sigma$). Rzeczywiście
$\operatorname{dist}(f(\sigma),\sigma)>3\delta/4$
(odległość dowolnego $f(x)$ od dowolnego $y\in\sigma$
wynosi co najmniej $\delta-\|x-y\|$).
Niech $K$ oznacza ustalony drobny podział obu kopii
wielościanu $X$. Dobierzmy następny podział $K'$ dziedziny
i aproksymację symplicjalną $g:K'\to K$ odwzorowania $f$.
Lemat o aproksymacji symplicjalnej zapewnia homotopię
$|g|\simeq f$ oraz, przy wystarczająco drobnej triangulacji
\emph{celu}, oszacowanie
$\sup_{x\in X}\||g|(x)-f(x)\|<\delta/4$.
Istotne jest rozróżnienie dwóch podziałów: $K'$ służy
do przybliżania odwzorowania, a drobny podział $K$ kontroluje
zarówno średnice sympleksów w bazie grup łańcuchów,
jak i błąd aproksymacji.

Niech $S:C_*(K;\Q)\to C_*(K';\Q)$ oznacza odwzorowanie
subdywizji. Złożenie $F=g_\#\circ S$ jest
odwzorowaniem łańcuchowym i indukuje w homologii
$f_*$. Każdy sympleks występujący w $S(\sigma)$
leży geometrycznie w $\sigma$. Jego obraz przez $g$
leży w odległości co najwyżej $\delta/4$ od
$f(\sigma)$, a zatem nie może być sympleksem $\sigma$,
ponieważ $f(\sigma)$ jest oddalone od $\sigma$
o ponad $3\delta/4$. W macierzy $F_k$ względem
bazy zorientowanych $k$-sympleksów wszystkie wyrazy
diagonalne są więc zerowe. Z Lematu~\ref{lem:hopf-trace-v10}
wynika $L(f)=\sum_k(-1)^k\operatorname{tr}F_k=0$,
wbrew założeniu. To dowodzi tezy dla wielościanów.
Pozostaje przypadek skończonego kompleksu CW $X$.
Korzystamy z klasycznego twierdzenia o otoczeniowych
retraktach. Przestrzeń metryzowalna $X$ jest \emph{absolutnym
retraktorem otoczeniowym} (ANR), gdy dla każdego
osadzenia $X$ jako domkniętego podzbioru przestrzeni
metryzowalnej $Y$ jest retraktem pewnego otwartego
otoczenia w $Y$. Skończony CW-kompleks jest zwartą
metryzowalną przestrzenią typu ANR; można go więc osadzić
jako domknięty podzbiór $i:X\hookrightarrow\R^N$
tak, że istnieje otwarte otoczenie $U$ i retrakcja
$r:U\to X$, $r\circ i=\id_X$. To jest osobne twierdzenie
o ANR, niewynikające z samego twierdzenia o aproksymacji
symplicjalnej.
Zwartość pozwala znaleźć skończony wielościan
$P\subset U$ zawierający $i(X)$: bierzemy skończenie
wiele dostatecznie małych sympleksów triangulacji
$\R^N$, których suma zawiera $i(X)$ i leży w $U$.
Niech $j:X\to P$ będzie tym osadzeniem, a
$R=r|_P:P\to X$ retrakcją. Definiujemy
$F=j\circ f\circ R:P\to P$.
Gdyby $F(y)=y$, to $y\in j(X)$ i z $R\circ j=\id_X$
wynikałby punkt stały $f$, sprzeczność.
Tak więc $F$ nie ma punktów stałych i z przypadku
wielościanowego $L(F)=0$.
Dla każdego $k$ zachodzi $F_*=j_* f_*R_*$,
a z cykliczności śladu także dla map między różnymi
skończeniewymiarowymi przestrzeniami
$\operatorname{tr}(j_*f_*R_*)=
\operatorname{tr}(f_*R_*j_*)=\operatorname{tr}(f_*)$.
Zatem $L(f)=L(F)=0$, co kończy dowód.
\end{proof}

\begin{wniosek}[Twierdzenie Brouwera]
Każde ciągłe odwzorowanie $f:D^n\to D^n$, $n\ge1$,
ma punkt stały.
\end{wniosek}
\begin{proof}
Dysk $D^n$ jest skończonym wielościanem z dokładnością
do homeomorfizmu i jest kontraktylny, więc
$L(f)=1$. Stosujemy Twierdzenie~\ref{thm:lefschetz-v10}.
\end{proof}

\begin{uwaga}[Właściwy zakres wniosku]
Odwrotność twierdzenia Lefschetza jest fałszywa:
$L(f)=0$ nie oznacza istnienia odwzorowania bez punktów
stałych w danej klasie homotopii ani nie oznacza, że
$f$ nie ma punktów stałych. Przykładowo identyczność
na okręgu ma $L(\id_{S^1})=\chi(S^1)=0$, ale każdy
punkt jest stały. Obrót okręgu przez kąt niezerowy
ma także $L=0$ i nie ma punktów stałych.
\end{uwaga}

\section{Indeks izolowanego punktu stałego}
\begin{definicja}[Indeks lokalny]
\label{def:fixed-index-v10}
Niech $M^n$ będzie gładką rozmaitością bez brzegu,
$n\ge1$, $f:M\to M$ odwzorowaniem ciągłym oraz
$p\in M$ izolowanym punktem stałym. Wybierzmy mapę
$\varphi:U\to\R^n$ wokół $p$ i małą kulę $B$
o domknięciu zawartym w $\varphi(U)$ taką, że
$\varphi\circ f\circ\varphi^{-1}$ jest określona
na $\overline B$ i $p$ jest jedynym punktem stałym
w $\varphi^{-1}(\overline B)$. Dla $F=\varphi f\varphi^{-1}$
określamy
$\operatorname{ind}_p(f)=\deg(\id-F,B,0)$,
tj. stopień Brouwera odwzorowania
$x\mapsto x-F(x)$ względem zera; równoważnie jest to
stopień odwzorowania
$\partial B\to S^{n-1}$, $x\mapsto(x-F(x))/\|x-F(x)\|$.
\end{definicja}

\begin{stwierdzenie}[Właściwości indeksu]
Indeks jest niezależny od wystarczająco małej mapy i kuli,
addytywny przy rozdzieleniu izolowanych punktów stałych
oraz niezmienniczy przy homotopii, podczas której
na brzegu rozważanego otoczenia nie pojawiają się punkty
stałe. Jeśli $f$ jest różniczkowalne w $p$
i $\det(I-\D f_p)\ne0$, to
$\operatorname{ind}_p(f)=\operatorname{sgn}\det(I-\D f_p)$.
\end{stwierdzenie}
\begin{proof}
Niezależność od otoczenia, addytywność i niezmienniczość
są własnościami stopnia Brouwera, wynikającymi
z ekscyzji, względnej homologii par $(B,\partial B)$
i homotopijnej niezmienniczości stopnia.
Dla ostatniej równości w lokalnych współrzędnych
$F(x)=F(0)+\D F_0x+o(\|x\|)$ oraz $F(0)=0$.
Jeśli $A=I-\D F_0$ jest odwracalna, wybieramy
$\varepsilon>0$ tak, by na sferze $\|x\|=\varepsilon$
reszta $o(\|x\|)$ miała normę mniejszą niż
$\tfrac12\min_{\|v\|=1}\|Av\|\,\varepsilon$.
Homotopia $Ax-t(F(x)-\D F_0x)$, $0\le t\le1$,
nie ma zer na brzegu kuli. Stopień $A$ wynosi
$\operatorname{sgn}\det A$, co daje wzór.
Niezależność od współrzędnych dla dowolnego odwzorowania ciągłego
wymaga czegoś więcej niż sprzężenia różniczek $f$.
Jeśli $h$ jest przejściem między dwoma układami współrzędnych
z $h(0)=0$, to w małym otoczeniu $0$ zachodzi
$h(x)-h(F(x))=(\D h_0+E(x))(x-F(x))$,
gdzie $\|E(x)\|\to0$ dla $x\to0$. Wynika to ze wzoru
na przyrost $h$ wzdłuż odcinka łączącego $x$ z $F(x)$;
w dostatecznie małej kuli odcinek leży w dziedzinie $h$.
Na jej brzegu macierze $\D h_0+tE(x)$, $0\le t\le1$,
są odwracalne. Stopień pola różnicowego zmienia się przy
przejściu do nowych współrzędnych o dwa czynniki
$\operatorname{sgn}\det\D h_0$: od zmiany współrzędnych
dziedziny oraz od mnożenia różnicy przez $\D h_0$.
Ich iloczyn wynosi $1$, stąd indeks nie zależy od mapy.
Także niezależność od małej kuli wynika z ekscyzji,
ponieważ jedynym zerem $\id-F$ pozostaje punkt $0$.
\end{proof}

\section{Formuła Lefschetza--Hopfa przez przecięcie grafu z diagonalną}
\begin{twierdzenie}[Formuła indeksowa]
\label{thm:lefschetz-hopf-v10}
Niech $M^n$ będzie zamkniętą gładką rozmaitością,
$n\ge1$, a $f:M\to M$ gładkim odwzorowaniem,
którego wszystkie punkty stałe są nieosobliwe,
tzn. $\det(I-\D f_p)\ne0$ dla każdego $f(p)=p$.
Wówczas punktów stałych jest skończenie wiele i
\[
 L(f)=\sum_{p:f(p)=p}\operatorname{sgn}\det(I-\D f_p).
\]
\end{twierdzenie}
\begin{proof}
Punkt stały $p$ odpowiada przecięciu grafu
$\Gamma_f=\{(x,f(x)):x\in M\}$ z diagonalną
$\Delta=\{(x,x):x\in M\}$ w $M\times M$.
Warunek $\det(I-\D f_p)\ne0$ jest dokładnie warunkiem
transwersalności obu $n$-wymiarowych podrozmaitości
w rozmaitości $2n$-wymiarowej. Przecięcia są więc
izolowane, a ponieważ $M$ jest zwarta, jest ich
skończenie wiele. Lokalny znak przecięcia jest
$\operatorname{sgn}\det(I-\D f_p)$: normalny kierunek
do diagonali identyfikujemy z różnicą pierwszego
i drugiego wektora stycznego.

Wyjaśnijmy równość globalnej sumy przecięć ze śladem
bez przemilczenia znaków. Najpierw rozpatrzmy spójną zorientowaną rozmaitość $M$.
Jeśli $M$ jest niespójna, to $f$ może przenosić jej spójne
składowe w inne składowe. W rozkładzie
$H^*(M;\Q)=\bigoplus_{\alpha}H^*(M_\alpha;\Q)$
śladowi $f^*$ odpowiadają tylko bloki diagonalne, a więc
te składowe $M_\alpha$, dla których $f(M_\alpha)\subseteq M_\alpha$.
Pozostałe składowe nie zawierają punktów stałych i nie
wnoszą wkładu do śladu. Dlatego argument dla przypadku
spójnego wystarczy na składowych zachowywanych przez $f$.
Dla rozmaitości nieorientowalnej potrzebna będzie odrębna
wersja dualności z lokalnym systemem orientacji;
nie wynika ona z rachunku dla zwykłej klasy fundamentalnej.
Niech $a_{p,j}$ będzie bazą $H^p(M;\Q)$ i $b_{p,j}$
bazą $H^{n-p}(M;\Q)$ dualną w sensie
$\langle a_{p,i}\smile b_{p,j},[M]\rangle=\delta_{ij}$.
Przez $\pi_1,\pi_2:M\times M\to M$ oznaczmy rzutowania.
Klasa diagonalna ma postać
\[
 \operatorname{PD}_{M\times M}[\Delta]
 =\sum_{p=0}^{n}\sum_j(-1)^{n-p}
 \pi_1^*a_{p,j}\smile\pi_2^*b_{p,j}.
\]
Wzór wynika z własności
$\langle\operatorname{PD}[\Delta]\smile
 \pi_1^*\alpha\smile\pi_2^*\beta,[M\times M]\rangle
 =\langle\alpha\smile\beta,[M]\rangle$
i reguły znaków Künnetha.
Dla odwzorowania grafu $j_f:M\to M\times M$,
$j_f(x)=(x,f(x))$, otrzymujemy
\begin{align*}
 \langle j_f^*\operatorname{PD}[\Delta],[M]\rangle
 &=\sum_{p,j}(-1)^{n-p}
 \langle a_{p,j}\smile f^*b_{p,j},[M]\rangle\\
 &=\sum_{q=0}^n(-1)^q
 \operatorname{tr}(f^*\mid H^q(M;\Q))=L(f).
\end{align*}
W drugim przejściu podstawiliśmy $q=n-p$;
macierze $f^*$ i $f_*$ mają te same ślady nad $\Q$.
Lewa strona jest globalną algebraiczną liczbą
przecięć grafu z diagonalną, a transwersalność
pozwala policzyć ją jako sumę znaków lokalnych.
Dla nieorientowalnego $M$ końcowa równość również jest
prawdziwa. Wymaga jednak odpowiednika klasy diagonalnej
z lokalnymi współczynnikami i sparowania skręconej
klasy Thoma normalnej do $\Delta$ z klasą fundamentalną
współczynnikową. Tę wersję twierdzenia indeksowego
przyjmujemy tutaj jako odrębny standardowy wynik;
powyższy rachunek bazami i wzorem Künnetha dowodzi
formuły bezpośrednio w przypadku orientowalnym.
Lokalna definicja indeksu pozostaje poprawna również
bez globalnego wyboru orientacji $M$.
\end{proof}

\begin{center}
\begin{tikzpicture}[>=Stealth,font=\small,scale=.95]
\draw[->] (-0.2,0)--(5.6,0) node[below] {$x$};
\draw[->] (0,-0.2)--(0,4.5) node[left] {$y$};
\draw[very thick,manifoldblue] (0,0)--(4.1,4.1)
 node[above right] {$\Delta=\{y=x\}$};
\draw[very thick,manifoldred,domain=0:4.2,samples=80]
 plot (\x,{2.6-0.5*\x+0.1*\x*\x});
\node[manifoldred,anchor=west] at (4.4,2.6) {$\Gamma_f$};
\fill[manifoldgreen] (2.0,2.0) circle (1.8pt);
\draw[dashed,manifoldgreen] (2.0,0)--(2.0,2.0);
\node[anchor=north] at (2.0,-.12) {punkt stały};
\end{tikzpicture}
\end{center}
\begin{center}\small Schemat lokalny: punkty stałe są przecięciami
wykresu $y=f(x)$ z diagonalną. Rysunek ilustruje mechanizm,
nie jest wykresem dowolnego odwzorowania rozmaitości.\end{center}

\section{Obliczenia: sfery, torusy i kontrprzykłady}
\begin{przyklad}[Sfera i stopień]
Jeśli $f:S^n\to S^n$ jest ciągłe, $n\ge1$,
to $H_0(S^n;\Q)\cong\Q$ i $H_n(S^n;\Q)\cong\Q$,
a na drugiej grupie $f_*$ jest mnożeniem przez $\deg f$.
Wobec tego $L(f)=1+(-1)^n\deg f$.
Dla odwzorowania antypodalnego $a(x)=-x$ mamy
$\deg a=(-1)^{n+1}$, więc $L(a)=0$, zgodnie z tym,
że $a$ nie ma punktów stałych.
Z drugiej strony każde odwzorowanie $S^{2m}\to S^{2m}$
stopnia różnego od $-1$ musi mieć punkt stały.
\end{przyklad}

\begin{przyklad}[Torus i całkowita macierz]
Niech $A\in M_d(\Z)$ i $f_A:T^d=\R^d/\Z^d\to T^d$
będzie mapą $[x]\mapsto[Ax]$. Ponieważ
$H_k(T^d;\Q)\cong\bigwedge^k\Q^d$,
odwzorowanie $f_{A*}$ działa jako $\bigwedge^k A$.
Tożsamość
$\det(I-tA)=\sum_{k=0}^d(-t)^k
 \operatorname{tr}(\bigwedge^kA)$ wynika z działania
operatora na potęgach zewnętrznych: w bazie klinów
$e_{i_1}\wedge\cdots\wedge e_{i_k}$ współczynniki
diagonalne $\bigwedge^kA$ są głównymi minorami stopnia $k$
macierzy $A$, a suma tych minorów jest współczynnikiem
wielomianu $\det(I-tA)$ przy $(-t)^k$.
Zatem $L(f_A)=\det(I-A)$.
Wzory na potęgi zewnętrzne i minory operatora
zostały również wyprowadzone w
\href{https://github.com/dariusz-piekarz/mathematical_scipts/blob/biblioteka-matematyczna-iteracja-6/chapters/pl/differential_geometry_and_topology/03-algebra-tensorowa-i-zewnetrzna.tex}{rozdziale o algebrze zewnętrznej},
pod historycznymi etykietami
\texttt{thm:potega-zewn-operatora} oraz
\texttt{thm:minory-lambda}. Nie powtarzamy tamtego dowodu
własności konstrukcji, lecz wykorzystujemy wynik w rachunku Lefschetza.
Jeśli $\det(I-A)\ne0$, liczba rozwiązań $[Ax]=[x]$
wynosi $|\det(I-A)|$; wszystkie są nieosobliwe,
a ich indeksy mają znak $\operatorname{sgn}\det(I-A)$.
Formuła Lefschetza--Hopfa daje więc tożsamość
$L(f_A)=\det(I-A)$ również jako rachunek lokalny.
\end{przyklad}

\begin{przyklad}[Dwa punkty stałe o ujemnych indeksach na torusie]
\label{prz:lefschetz-torus-dwa-indeksy}
Niech $T^2=\R^2/\Z^2$ i rozważmy
odwzorowanie $f_A([x])=[Ax]$ wyznaczone przez
\[
 A=\begin{pmatrix}2&1\\1&0\end{pmatrix}.
\]
Ma ono wyznacznik $-1$, zatem jest dyfeomorfizmem
torusa odwracającym orientację. Liczbę Lefschetza
obliczamy bezpośrednio z homologii:
\[
 \operatorname{tr}(f_\ast\mid H_0)=1,\qquad
 \operatorname{tr}(f_\ast\mid H_1)=\operatorname{tr}A=2,
 \qquad
 \operatorname{tr}(f_\ast\mid H_2)=\det A=-1,
\]
skąd $L(f_A)=1-2-1=-2$.
Zgadza się to z wyznacznikiem
\[
 \det(I-A)
 =\det\begin{pmatrix}-1&-1\\-1&1\end{pmatrix}=-2.
\]
Sprawdźmy, że nie ukrywa się tu wielokrotność
jednego punktu stałego. Równanie $[Ax]=[x]$
jest równoważne $(A-I)x\in\Z^2$.
Macierz
$A-I=\bigl(\begin{smallmatrix}1&1\\1&-1\end{smallmatrix}\bigr)$
jest odwracalna nad $\R$. Jej obraz kratowy
$(A-I)\Z^2$ składa się dokładnie z par
liczb całkowitych o jednakowej parzystości:
dla $(m,n)\in\Z^2$ mamy
$(m+n,m-n)$, a odwrotnie z pary $(u,v)$
tej samej parzystości odzyskujemy
$m=(u+v)/2$ i $n=(u-v)/2$.
Iloraz $\Z^2/(A-I)\Z^2$ ma więc dwa elementy.
Ich reprezentantom odpowiadają dokładnie
\[
 [x_0]=[(0,0)],\qquad
 [x_1]=[(1/2,1/2)].
\]
Istotnie,
$(A-I)(1/2,1/2)=(1,0)\in\Z^2$.

W obu punktach stałych różniczka wynosi $A$,
a ponieważ $\det(I-A)=-2\ne0$, są one
nieosobliwe i każdy ma lokalny indeks $-1$.
Suma indeksów to $-1-1=-2=L(f_A)$,
zgodnie z Twierdzeniem~\ref{thm:lefschetz-hopf-v10}.
Przykład pokazuje, że dodatnia liczba punktów stałych
nie musi oznaczać dodatniej liczby Lefschetza.
\end{przyklad}

\begin{figure}[!htbp]
\centering
\begin{tikzpicture}[>=Stealth,font=\small]
 \draw[thick,fill=manifoldblue!5] (0,0) rectangle (4.4,4.4);
 \draw[gray,densely dashed] (2.2,0)--(2.2,4.4);
 \draw[gray,densely dashed] (0,2.2)--(4.4,2.2);
 \fill[manifoldred] (0,0) circle (3pt);
 \fill[manifoldred] (2.2,2.2) circle (3pt);
 \node[anchor=north west] at (.12,-.12) {$[(0,0)]$};
 \node[anchor=south west] at (2.32,2.3)
       {$[(\tfrac12,\tfrac12)]$};
 \node[below] at (2.2,-.65)
       {Kwadrat fundamentalny $\R^2/\Z^2$};
 \node[right] at (4.52,3.8)
       {$\operatorname{ind}=-1$};
 \node[right] at (4.52,3.2)
       {dla obu punktów};
\end{tikzpicture}
\caption{Dwa punkty stałe odwzorowania torusa
o macierzy $\left(\begin{smallmatrix}2&1\\1&0\end{smallmatrix}\right)$.
Punkty na przeciwległych krawędziach kwadratu
należy utożsamić. Oba lokalne indeksy są ujemne.}
\label{fig:lefschetz-torus-dwa-indeksy}
\end{figure}

\begin{przyklad}[Rotacja okręgu]
Dla rotacji $R_\alpha:S^1\to S^1$ z $\alpha\notin2\pi\Z$
odwzorowanie jest homotopijne do identyczności i ma
$L(R_\alpha)=1-1=0$; nie ma też punktu stałego.
Dla $\alpha=0$ ta sama liczba Lefschetza wynosi zero,
choć wszystkie punkty są stałe. Dowodzi to, że
warunek $L(f)\ne0$ jest wystarczający, lecz niekonieczny,
dla istnienia punktu stałego.
\end{przyklad}

\begin{uwaga}[Zakres dowodu formuły indeksowej]
Dowód formuły Lefschetza--Hopfa w przypadku zorientowanym
jest redukcją do dualności Poincarégo i klasy diagonalnej,
nie nowym dowodem dualności. W przypadku nieorientowalnym
korzystamy z wersji indeksowej uwzględniającej lokalne
współczynniki. Dla skończonych CW-kompleksów samą
implikację $L(f)\ne0\Rightarrow\operatorname{Fix}(f)\ne\varnothing$
udowodniliśmy niezależnie, stosując aproksymację
symplicjalną i własność ANR.
\end{uwaga}

\begin{przyklad}[Lokalny indeks a liczba Lefschetza: odbicie odcinka]
Niech $f:[-1,1]\to[-1,1]$, $f(x)=-x$. Jedynym
punktem stałym jest $0$, a dla pola różnicowego
$\id-f=2x$ jego stopień lokalny jest równy $+1$.
Z drugiej strony $L(f)=1$, ponieważ przedział jest
ściągalny. Przykład ilustruje zgodność obu metod
w przestrzeni z brzegiem, lecz nie stosujemy do niego
bezpośrednio formuły dla zamkniętej rozmaitości:
wyznaczamy indeks w otwartym wnętrzu i używamy
osobno twierdzenia Lefschetza dla skończonych kompleksów.
\end{przyklad}

\section{Relacje z teorią Morse'a i dalszymi częściami książki}
Dla funkcji Morse'a liczba $\chi(M)=\sum_p(-1)^{\lambda(p)}$
jest naprzemienną sumą generatorów kompleksu Morse'a.
Z drugiej strony $L(\id_M)=\chi(M)$ jest śladem działania
identyczności na homologii. To dwa geometrycznie odmienne
wyprowadzenia tego samego niezmiennika. Połączeniem obu
jest twierdzenie Poincarégo--Hopfa: suma lokalnych indeksów
zer pola wektorowego jest równa $\chi(M)$.
Nie utożsamiamy tego twierdzenia z przypadkiem odwzorowania
$\id_M$ w formule indeksowej, gdyż identyczność nie
ma izolowanych punktów stałych. Należałoby najpierw
wybrać odpowiednie zaburzenie identyczności i porównać
lokalne stopnie; taki argument należy do teorii indeksu
pól wektorowych.

% Źródła orientacyjne: A. Hatcher, Algebraic Topology;
% J. Milnor, Morse Theory; M. Schwarz, Morse Homology;
% R. Brown, Lefschetz Fixed Point Theorem; D. Kim, The Lefschetz fixed-point formula.


% Iteracja XI: szczegółowy audyt części analitycznej homologii Morse'a.
